%% 33장 수치 미분
\chapter{수치 미분}
\label{ch:numdiff}

\abstract{함수의 값만 계산할 수 있을 때 도함수를 어림하는 가장 쉬운 방법은 가까운 두 점의 기울기를 재는 것이다. 그러나 간격을 줄이면 절단 오차는 줄지만 두 값의 차에서 생기는 소거가 반올림 오차를 키워 어떤 간격 아래로는 결과가 오히려 나빠진다. 이 장에서는 임의의 점 배치에서 차분 공식의 무게를 보간 다항식의 도함수로 유도하고, 절단 오차와 반올림 오차의 균형에서 최적의 간격과 도달할 수 있는 정확도가 $\varepsilon^{p/(p+1)}$로 정해진다는 것을 보인다. 간격을 줄여 가며 오차의 선행 항을 없애는 Richardson 외삽, 소거가 아예 없는 복소 간격 미분, 간격을 자료에 맞게 고르는 방법을 다루고, stardust가 차분 공식의 무게를 정확한 유리수로 계산해 한 번만 반올림하는 이유를 본다.}

\section{차분 공식}
\label{sec:fd-formulas}

\subsection{앞쪽 차분과 중앙 차분}

도함수의 정의를 그대로 쓰면 \term{앞쪽 차분}{forward difference}
\[
D_+(h)=\frac{f(x+h)-f(x)}{h}=f'(x)+\frac h2f''(x)+O(h^2)
\]
을 얻는다. 오차는 Taylor 전개에서 바로 나오며 $h$에 비례한다. 양쪽으로 같은 간격을 두면 짝수 차수의 항이 상쇄되어
\[
D_0(h)=\frac{f(x+h)-f(x-h)}{2h}=f'(x)+\frac{h^2}6f'''(x)+\frac{h^4}{120}f^{(5)}(x)+\cdots
\]
이 된다. 이것이 \term{중앙 차분}{central difference}이고, 오차가 $h^2$에 비례하므로 같은 간격에서 훨씬 정확하다. 차분 공식의 오차를 $h^p$로 쓸 때 $p$를 공식의 차수라 한다. 중앙 차분의 오차 전개에는 짝수 거듭제곱만 있는데 뒤의 Richardson 외삽은 이 성질을 이용한다.

\subsection{임의의 스텐실}

더 많은 점을 쓰면 차수를 더 올릴 수 있다. $x+t_ih$ $(i=0,\dots,n)$에서 계산한 값으로 $f'(x)$를 어림하는 공식 $h^{-1}\sum_iw_if(x+t_ih)$를 생각하자. 점들의 상대 위치 $t_i$를 \term{스텐실}{stencil}이라 한다. 가장 자연스러운 무게는 29장의 보간 다항식을 미분해서 얻는다.

\begin{proposition}
\label{prop:fd-weights}
서로 다른 $t_0,\dots,t_n$의 Lagrange 기저 다항식을 $L_i$라 하고 $w_i=L_i'(0)$으로 두면, 공식 $h^{-1}\sum_iw_if(x+t_ih)$는 차수 $n$ 이하의 다항식에서 정확하고, $f\in C^{n+1}$이면 오차가 $O(h^n)$이다. 또 이 $w$는 Vandermonde 계
\[
\sum_{i=0}^nw_it_i^k=\delta_{k1},\qquad k=0,\dots,n
\]
의 유일한 해이다.
\end{proposition}

\begin{proof}
$g(t)=f(x+th)$의 보간 다항식은 $p(t)=\sum_ig(t_i)L_i(t)$이고, $p'(0)=\sum_iw_ig(t_i)$이다. $g'(0)=hf'(x)$이므로 공식은 $p'(0)/h$이다. $f$가 차수 $n$ 이하의 다항식이면 $g$도 그렇고 $p=g$이므로 정확하다. 단항식 $t^k$에 적용하면 $\sum_iw_it_i^k=(t^k)'\big|_{t=0}=\delta_{k1}$이므로 $w$는 이 계를 만족하고, 계의 행렬은 정칙인 Vandermonde 행렬의 전치이므로 해는 유일하다. 오차는 다음과 같이 얻는다. $e=g-p$는 $t_0,\dots,t_n$에서 $0$이므로 Rolle 정리에 의해 $e'$은 그 사이에 근 $\eta_1,\dots,\eta_n$을 가진다. 그러면 차수 $n-1$ 이하의 $p'$은 노드 $\eta_j$에서 $g'$을 보간하는 다항식이므로, 29장의 \cref{thm:interp-error}를 $g'$과 노드 $\eta_j$에 쓰면 $e'(0)=g^{(n+1)}(\zeta)\prod_j(0-\eta_j)/n!$인 $\zeta$가 있다. $|\eta_j|$가 스텐실의 폭으로 묶이고 $g^{(n+1)}=h^{n+1}f^{(n+1)}$이므로, $h$로 나누면 $O(h^n)$이다.
\end{proof}

스텐실이 원점에 대해 대칭이면 무게가 반대칭이 되어, 짝수 차수의 단항식에서는 합 $\sum_iw_it_i^k$가 저절로 $0$이다. 그래서 점이 짝수 개, 곧 $n$이 홀수이면 $t^{n+1}$에서도 정확해져 차수가 하나 더 오른다. 두 점 $\pm1$을 쓰는 중앙 차분이 $n=1$인데 차수가 $2$인 것이 그 경우이다. 가운데 점 $0$을 더해도 그 무게가 $0$이어서 차수는 그대로이다. 무게를 Vandermonde 계로 구하는 것은 29장에서 본 대로 조건수가 나쁜 방법이다. Fornberg가 내놓은 안정한 알고리즘은 노드를 하나씩 더해 가며 모든 차수의 도함수 무게를 점화식으로 만든다\cite{fornberg1988}. 스텐실이 고정되어 있으면 정확한 유리수 산술로 계산하는 것이 가장 확실하다.

\section{절단 오차와 반올림 오차}
\label{sec:fd-rounding}

\subsection{두 오차의 균형}

차분 공식의 계산값에는 절단 오차 말고도 반올림 오차가 있다. $f$를 계산한 값이 $\hat f=f(1+\delta)$, $|\delta|\le\eta$라 하자. 정확히 반올림한 기본 함수라면 $\eta\approx u$이고, 복잡한 함수라면 더 크다. 앞쪽 차분의 분자는 거의 같은 두 값의 차이므로 9장의 소거가 일어나고 반올림 오차가 $h$로 나뉘어 커진다. 오차 전체는 대략
\begin{equation}
\label{eq:fd-total-error}
E(h)\approx\frac h2|f''|+\frac{2\eta|f|}h
\end{equation}
이고, 첫 항은 $h$와 함께 줄지만 둘째 항은 $h$가 작아질수록 커진다. 두 항의 합을 최소로 하는 간격은 $h^\ast=2\sqrt{\eta|f|/|f''|}$이고 그때의 오차는 $2\sqrt{\eta|f||f''|}$이다. $f$와 $f''$의 크기가 비슷하면 앞쪽 차분은 $h\approx\sqrt\eta$에서 많아야 $\sqrt\eta$의 상대오차에 이른다. 일반적으로 차수 $p$의 공식은 절단 오차가 $Ch^p$, 반올림 오차가 $\eta/h$에 비례하므로
\[
h^\ast\sim\eta^{1/(p+1)},\qquad E(h^\ast)\sim\eta^{p/(p+1)}
\]
이다. 중앙 차분($p=2$)은 $h\approx\varepsilon^{1/3}\approx6\times10^{-6}$에서 $\varepsilon^{2/3}\approx4\times10^{-11}$ 정도의 정확도를 낸다. 차수를 올릴수록 도달할 수 있는 정확도는 $\varepsilon$에 다가가지만 결코 $\varepsilon$에 이르지는 못한다. 반올림 오차의 항에 언제나 $1/h$가 붙기 때문이다.

$x=1$에서 $e^x$의 도함수를 $h=10^{-k}$로 어림하면 이 균형이 그대로 보인다. 앞쪽 차분의 상대오차는 $k=1$에서 $8$까지 $5.2\times10^{-2}$, $5.0\times10^{-3}$, $\dots$, $5.1\times10^{-8}$, $2.4\times10^{-9}$으로 $h$에 비례해 줄다가, $k=9$부터 $7.9\times10^{-8}$, $5.7\times10^{-7}$, $1.2\times10^{-5}$로 다시 커지고 $k=15$에서는 $0.14$이다. 가장 좋은 값은 $h=10^{-8}\approx\sqrt\varepsilon$ 근처의 $2.4\times10^{-9}$이다. 중앙 차분은 $k=5$에서 $2.2\times10^{-11}$로 가장 좋고 그 뒤로는 반올림이 이긴다. 두 수치 모두 위의 예측 $\sqrt\varepsilon\approx1.5\times10^{-8}$, $\varepsilon^{2/3}\approx3.7\times10^{-11}$과 맞는다.

\subsection{복소 간격 미분}

$f$가 실수축 근처에서 해석적이고 복소수 인자로 계산할 수 있으면 소거를 아예 피할 수 있다. Taylor 전개 $f(x+ih)=f(x)+ihf'(x)-h^2f''(x)/2-ih^3f'''(x)/6+\cdots$의 허수부를 취하면
\[
f'(x)=\frac{\operatorname{Im}f(x+ih)}h+O(h^2)
\]
이다. 뺄셈이 없으므로 반올림 오차가 $h$로 나뉘어 커지지 않고, $h$를 얼마든지 작게 잡아 절단 오차를 없앨 수 있다. 같은 $e^x$에서 $h=10^{-10}$이면 상대오차가 $1.6\times10^{-16}$이고, $h=10^{-20}$이나 $10^{-100}$에서는 정확히 맞는다. 이 \term{복소 간격 미분}{complex-step derivative}은 Lyness와 Moler가 해석 함수의 수치 미분으로 처음 다루었고\cite{lyness1967}, Squire와 Trapp이 실용적인 방법으로 널리 알렸다\cite{squire1998,martins2003}. 다만 함수의 모든 연산이 복소수를 받아야 하고 절댓값이나 비교처럼 해석적이지 않은 연산이 섞이면 틀린다. 일반적인 라이브러리 함수가 그 조건을 만족하리라 기대하기는 어렵다.

\section{Richardson 외삽}
\label{sec:richardson}

중앙 차분의 오차 전개가 $D_0(h)=f'(x)+c_1h^2+c_2h^4+\cdots$이면, 두 간격의 값을 결합해 선행 항을 없앨 수 있다.
\[
D_1(h)=D_0(h/2)+\frac{D_0(h/2)-D_0(h)}{4-1}=f'(x)+O(h^4)
\]
이것을 되풀이해 $j$번째 열에서는 $D_j(h)=D_{j-1}(h/2)+\bigl(D_{j-1}(h/2)-D_{j-1}(h)\bigr)/(4^j-1)$로 $h^{2j}$의 항을 없애는 것이 \term{Richardson 외삽}{Richardson extrapolation}이다\cite{richardson1911,richardson1927}. 오차가 $h$의 거듭제곱 급수로 전개되는 모든 계산에 쓸 수 있으며 수치 적분의 Romberg 방법도 같은 생각이다.

$x=1$에서 $e^x$의 중앙 차분을 $h=0.4$에서 시작해 반씩 줄이며 표를 채우면 상대오차가 다음과 같다. 첫 열은 $2.7\times10^{-2}$, $6.7\times10^{-3}$, $1.7\times10^{-3}$, $4.2\times10^{-4}$, $1.0\times10^{-4}$로 네 배씩 줄지만 대각선을 따라가면 $5.4\times10^{-5}$, $1.3\times10^{-8}$, $4.4\times10^{-13}$, $1.8\times10^{-15}$로 내려간다. 가장 작은 간격이 $0.025$에 불과한데도 기계 정밀도 근처에 이른다. 외삽의 진짜 이점은 간격을 크게 둘 수 있어 반올림 오차가 작다는 데 있다. 또 표의 이웃한 값의 차이가 오차의 좋은 추정이 되므로 값이 더 이상 좋아지지 않는 곳에서 멈출 수 있다. Ridders는 이 표를 적응적으로 채우다가 오차 추정이 커지기 시작하면 멈추자고 제안했다\cite{ridders1982}.

\section{간격 고르기}
\label{sec:step-choice}

\subsection{척도}

\eqref{eq:fd-total-error}의 최적 간격은 $|f|/|f''|$, 곧 함수가 변하는 길이의 척도에 달려 있다. 이 척도는 보통 알 수 없으므로 실제 라이브러리는 대개 $x$의 크기를 척도로 쓴다. scipy의 \code{least\_squares}가 쓰는 \code{approx\_derivative}의 기본 간격은 앞쪽 차분에서 $h=\sqrt\varepsilon\,\max(1,|x|)$, 3점 중앙 차분에서 $\varepsilon^{1/3}\max(1,|x|)$이다. \code{approx\_fprime}은 척도 없이 절대 간격 $\sqrt\varepsilon$을 쓴다. 이 기본값은 대부분의 경우 무난하지만 $x$의 크기가 함수의 척도와 다를 때 실패한다. 예컨대 $x=10^6$에서 $\sin(1000x)$를 미분하면 $h=\sqrt\varepsilon\cdot10^6\approx1.5\times10^{-2}$인데, 이것은 함수의 한 주기 $2\pi/1000\approx6.3\times10^{-3}$보다 크다. 거꾸로 $x$가 $0$에 아주 가깝고 함수가 $x$의 크기와 무관한 척도로 변하면, $|x|$에 비례하는 간격은 지나치게 작아진다.

간격의 표현 가능성도 신경 써야 한다. $x+h$가 binary64로 반올림되므로 실제로 쓰인 간격은 $h$가 아니라 $(x+h)-x$이다. 간격을 2의 거듭제곱으로 잡고 $x$의 지수에 맞추면 이 문제가 없다. Gill, Murray, Saunders, Wright는 최적화에서 함수의 2차 차분을 몇 번 계산해 \eqref{eq:fd-total-error}의 균형점을 직접 찾자고 제안했다\cite{gill1983}.

\subsection{적응적 간격}

scipy의 \code{derivative}는 Richardson 외삽 대신 높은 차수의 공식을 간격을 줄여 가며 되풀이한다. 기본은 $8$차 중앙 공식이고, 스텐실은 $x\pm h/c^k$ $(k=0,\dots,3)$, 간격은 반복마다 $c=2$배씩 줄인다. 이렇게 하면 다음 반복에서 새로 계산할 점이 둘뿐이다. 이웃한 두 반복의 차이를 오차의 추정으로 삼아 그 추정이 허용오차 아래로 내려가면 멈추고 직전의 열 배를 넘으면 반올림이 절단 오차를 이기기 시작한 것으로 보고 멈춘다.

이 스텐실의 노드는 $0,\pm1,\pm1/2,\pm1/4,\pm1/8$이다. \cref{prop:fd-weights}의 무게를 numpy로 Vandermonde 계를 풀어 구하면 정확한 유리수 무게와 아홉 개 모두 다르고 최대 상대오차가 $2.9\times10^{-13}$이다. 가운데 무게는 정확히 $0$이어야 하는데 $-3.3\times10^{-16}$이 나온다. $12$차 공식이면 상대오차가 $5.7\times10^{-8}$까지 커진다. 무게의 오차는 결과의 상대오차에 그대로 더해지므로 높은 차수의 공식으로 얻으려던 정확도의 일부를 무게를 계산하는 단계에서 잃는다.

\begin{instardust}[\code{docs/design/integrate.md} §3.8]
stardust의 \code{sd.differentiate.derivative}는 scipy의 알고리즘과 멈춤 판정을 원소마다 옮겼다. 다만 무게는 Vandermonde 계를 풀지 않고 \cref{prop:fd-weights}의 $L_i'(0)$을 정확한 유리수로 계산해 한 번 반올림한다. 중앙 스텐실의 가운데 무게는 정확히 $0$이고 나머지는 정확히 반대칭이다. 한쪽 스텐실은 노드가 $\sqrt c$의 거듭제곱이므로 $\mathbb Q(\sqrt c)$에서 정확히 계산해 올바르게 반올림한다. 무게는 $(n,c)$마다 한 번 계산해 기억한다. 짝을 이룬 두 점의 차 $w\bigl(f(x-ht)-f(x+ht)\bigr)$를 바깥 짝부터 더하므로 원소의 결과는 배치의 크기나 스레드 수와 무관하다. 설계 문서에 적힌 부수 효과 하나는 $0$에서 $x^3$의 도함수를 구할 때 나타난다. 정확한 무게로는 반복마다 같은 작은 값이 나와 오차 추정이 $0$이 되어 둘째 반복에서 수렴하지만, scipy는 무게의 반올림 흔들림 때문에 기본 허용오차로 끝내 멈추지 못한다.
\end{instardust}

\begin{instardust}[\code{docs/design/optimize.md}]
최적화와 비선형 회귀가 쓰는 야코비안은 다른 방법이다. 5점 중앙 차분이며 간격은 변수마다 처음에 한 번 고른다. 그 간격은 2의 거듭제곱이고 $2^{-8}|x_j|$에서 시작해 네 배씩 줄이며 이웃한 두 추정의 차이가 직전의 $1/8$보다 작아지는 동안 계속한다. 처음 열이 반올림 수준이면, 곧 $x_j$가 함수가 변하는 척도보다 훨씬 작으면, 간격을 키우는 쪽으로 다시 찾는다. 이렇게 고른 간격의 \code{approx\_fprime}은 상대오차가 $10^{-11}$--$10^{-13}$이다. scipy의 앞쪽 차분은 변수마다 함수를 한 번만 부르므로 약 아홉 배 빠르지만 상대오차가 약 $10^{-8}$이어서, 설계 문서는 정확한 쪽을 기본으로 두고 앞쪽 차분은 \code{epsilon=2**-26}으로 고를 수 있게 했다.
\end{instardust}

\section{야코비안과 헤세 행렬}
\label{sec:jacobian}

$F:\mathbb R^n\to\mathbb R^m$의 야코비안은 열마다 한 변수의 차분이므로 앞쪽 차분이면 함수를 $n+1$번, 중앙 차분이면 $2n$번 부른다. 야코비안이 희소하면 서로 다른 행에만 원소가 있는 열들은 한꺼번에 흔들어도 된다. 예컨대 삼중대각 야코비안은 열 $j$, $j+3$, $j+6$, $\dots$을 같은 차분으로 재면 셋의 결과가 서로 다른 행에 나타난다. 이렇게 열을 묶는 방법은 Curtis, Powell, Reid가 처음 제안했고\cite{curtis1974}, Coleman과 Moré가 그것이 그래프 색칠 문제임을 보였다\cite{coleman1983}. 헤세 행렬은 기울기의 야코비안이므로 차분을 두 번 겹쳐 얻는다. 반올림 오차가 두 번 나뉘므로 같은 방식으로는 흔히 $\varepsilon^{1/2}$이나 그보다 나쁜 정확도밖에 얻지 못한다.

함수의 계산 과정을 알 수 있다면 차분 대신 \term{자동 미분}{automatic differentiation}이 정확한 도함수를 기계 정밀도로 준다. 연쇄 법칙을 연산 하나하나에 적용하는 것이다. 이 라이브러리는 사용자가 넘긴 Python 함수를 상자로 보므로 자동 미분을 하지 않지만 그 이론은 Griewank와 Walther의 책\cite{griewank2008}이 자세하다.

\section{맺음}

수치 미분은 수치해석에서 가장 단순한 계산이지만 반올림 오차도 가장 뚜렷하게 드러난다. 간격은 줄인다고 늘 좋아지지 않는다. 절단 오차와 반올림 오차가 균형을 이루는 곳에 최적의 간격이 있고 그 균형이 도달할 수 있는 정확도를 $\varepsilon^{p/(p+1)}$로 묶는다. Richardson 외삽과 높은 차수의 공식은 큰 간격으로도 높은 정확도를 얻어 이 한계를 뒤로 미루고 복소 간격 미분은 소거를 없애 한계에서 아예 벗어난다. 다음 장에서는 도함수를 써서 방정식의 근을 찾는다. Newton 방법에 도함수가 없으면 여기서 본 차분이 그 자리를 대신한다.

수치 미분의 실제적인 측면은 Dennis와 Schnabel\cite{dennis1996}, Nocedal과 Wright\cite{nocedal2006}의 최적화 책이 잘 다룬다.

\chapterreferences
